\clearpage
\phantomsection% 
\addcontentsline{toc}{chapter}{KATA PENGANTAR}
%\thispagestyle{fancy}

\begin{justifying}
	\large \bfseries \centering \MakeUppercase{Kata Pengantar}\linebreak
	
	\normalsize \normalfont \justifying
	Puji dan syukur penulis panjatkan kepada Tuhan Yang Maha Esa atas segala berkat, rahmat, dan penyertaan-Nya sehingga penyusunan tugas akhir yang berjudul ``Pengembangan Sistem Terintegrasi Penerimaan Mahasiswa Baru Berbasis Web Menggunakan Metode Prototype (Studi Kasus: Universitas HKBP Nommensen)'' dapat diselesaikan dengan baik. \par
	Tugas akhir ini disusun sebagai salah satu syarat untuk memperoleh gelar Sarjana Komputer pada Program Studi Teknik Informatika, Fakultas Teknologi Industri, Institut Teknologi Sumatera. Dalam proses penyusunan tugas akhir ini, penulis telah mendapatkan banyak bantuan, bimbingan, dan dukungan dari berbagai pihak. \par
	Oleh karena itu, pada kesempatan ini penulis mengucapkan terima kasih kepada: \par
	\begin{enumerate}
		\item Tuhan Yang Maha Esa, atas segala berkat dan penyertaan-Nya selama penyusunan tugas akhir ini.
		\item Orang tua dan keluarga tercinta, atas doa, dukungan moral, dan motivasi yang tiada henti selama menempuh pendidikan.
		\item Bapak/Ibu Dosen Pembimbing, atas bimbingan, arahan, masukan, dan waktu yang telah diberikan selama penyusunan tugas akhir ini.
		\item Dosen dan staf pengajar di Program Studi Teknik Informatika Institut Teknologi Sumatera, atas ilmu pengetahuan dan pengalaman yang telah diberikan selama masa perkuliahan.
		\item Pihak Universitas HKBP Nommensen yang telah memberikan izin dan kesempatan untuk melakukan penelitian di lingkungan kampus.
		\item Teman-teman seperjuangan dan seluruh civitas akademika Institut Teknologi Sumatera yang telah memberikan semangat dan bantuan selama proses pengerjaan tugas akhir.
	\end{enumerate} \par
	Penulis menyadari bahwa tugas akhir ini masih jauh dari sempurna. Oleh karena itu, kritik dan saran yang membangun sangat diharapkan untuk perbaikan di masa mendatang. Semoga tugas akhir ini dapat memberikan manfaat dan kontribusi positif bagi perkembangan ilmu pengetahuan dan teknologi, khususnya di bidang Teknik Informatika.
	\vfill
	
\end{justifying}
\clearpage
